\documentclass[10pt]{extarticle}
\usepackage[a4paper,margin=1.5cm]{geometry}
\usepackage[T1]{fontenc}
\usepackage{amsmath,amssymb}
\usepackage{graphicx}
\usepackage[export]{adjustbox}
\usepackage{multicol}
\usepackage{enumitem}
\usepackage{booktabs,array,tabularx}
\setlength{\columnsep}{0.7cm}
\setlength{\parindent}{0pt}
\setlength{\parskip}{3pt}
\pagestyle{empty}
\makeatletter
% figure registry (figdefs.tex is generated by the extract step)
\newcommand{\qfigdef}[4]{\@namedef{qf@file@#1}{#2}\@namedef{qf@w@#1}{#3}\@namedef{qf@wide@#1}{#4}}
\newcommand{\qf@check}[1]{\@ifundefined{qf@file@#1}{\PackageError{score2}{Unknown figure id #1}{}}{}}
\newcommand{\qf@img}[4][]{\includegraphics[width=#2\dimexpr\@nameuse{qf@w@#4}pt\relax,max width=#3,max height=0.45\textheight,#1]{\@nameuse{qf@file@#4}}}
% Full-width questions: finalize wraps every question that contains a wide figure in \qfullbegin...\qfullend,
% so the WHOLE question (stem, figures, options) is set across the page in reading order, then the two columns
% resume. The question's head - its number, stem and first figure (\qfullsplit marks the end of the head) -
% is kept on one page, so a number is never left at a page foot with its figure overleaf; the rest may break
% across pages normally, so no big gaps are left. Without a \qfullsplit, the whole question is kept together
% if it is under 45% of a page. Never write these yourself.
\newif\ifqf@full
\newif\ifqf@boxed
\newsavebox{\qf@box}
\newcommand{\qf@open}{\global\qf@fulltrue\par\global\qf@boxedtrue
  \setbox\qf@box\vbox\bgroup\hsize\textwidth\linewidth\textwidth}
\newcommand{\qf@emit}[1]{\par\egroup\global\qf@boxedfalse
  \ifdim\dimexpr\ht\qf@box+\dp\qf@box\relax<#1\textheight\noindent\usebox{\qf@box}\par\else\unvbox\qf@box\fi}
\newcommand{\qfullsplit}{\ifqf@boxed\qf@emit{0.85}\fi}
\newcommand{\qf@close}{\ifqf@boxed\qf@emit{0.45}\fi\global\qf@fullfalse}
\newcommand{\qfullbegin}{\end{multicols}\qf@open}
\newcommand{\qfullend}{\qf@close\begin{multicols}{2}}
\newcommand{\qfullstart}{\qf@open}
\newcommand{\qfulldone}{\qf@close}
% \qfig[scale]{id}: block figure, centred, directly where it is written (size set by the tool).
% Fallback if a wide figure is not inside a full-width question: break out of the columns in place - no floats.
\newcommand{\qfig}[2][1]{\qf@check{#2}%
  \ifqf@full\qf@blk{#1}{#2}\else
  \ifnum\@nameuse{qf@wide@#2}=1 \expandafter\@firstoftwo\else\expandafter\@secondoftwo\fi
  {\end{multicols}\begin{center}\qf@img{#1}{\textwidth}{#2}\end{center}\begin{multicols}{2}}%
  {\qf@blk{#1}{#2}}\fi}
\newcommand{\qf@blk}[2]{\par\smallskip{\centering\qf@img{#1}{\linewidth}{#2}\par}\smallskip}
% \qfigi[scale]{id}: inline figure (inside option lists / tables / minipages); never breaks columns.
\newcommand{\qfigi}[2][1]{\qf@check{#2}\qf@img[valign=c]{#1}{\linewidth}{#2}}
\makeatother
% structure macros - lint reads these to check completeness and figure placement
\newcommand{\papersec}[1]{\par\bigskip{\centering\Large\bfseries PAPER-#1\par}\medskip}
\newcommand{\partsec}[2]{\par\medskip\noindent{\large\bfseries PART-#1: #2}\par\smallskip}
\newcommand{\sect}[1]{\par\smallskip\noindent\textbf{#1}\par}
\newcommand{\qn}[1]{\par\smallskip\noindent\textbf{#1.}~}
\newcommand{\paraq}[2]{\par\smallskip\noindent\textbf{Paragraph for Questions #1 and #2}\par}
\newenvironment{qoptions}{\begin{itemize}[leftmargin=1.8em,itemsep=1pt,topsep=2pt,label={}]}{\end{itemize}}
\newcommand{\opt}[2]{\item[(#1)] #2}
% figure registry (generated by score2_tools.py finalize)
\qfigdef{p1_phy_s1_q02_1}{Score Advanced - Test 03 (27 Jun 2021)-img/p1_phy_s1_q02_1.png}{75.8}{0}
\qfigdef{p1_phy_s1_q04_1}{Score Advanced - Test 03 (27 Jun 2021)-img/p1_phy_s1_q04_1.png}{131.2}{0}
\qfigdef{p1_phy_s1_q06_1}{Score Advanced - Test 03 (27 Jun 2021)-img/p1_phy_s1_q06_1.png}{59.9}{0}
\qfigdef{p1_phy_s2_para8_10_1}{Score Advanced - Test 03 (27 Jun 2021)-img/p1_phy_s2_para8_10_1.png}{95.6}{0}
\qfigdef{p1_phy_s2_para8_10_2}{Score Advanced - Test 03 (27 Jun 2021)-img/p1_phy_s2_para8_10_2.png}{94.5}{0}
\qfigdef{p1_phy_s2_para8_10_3}{Score Advanced - Test 03 (27 Jun 2021)-img/p1_phy_s2_para8_10_3.png}{93.7}{0}
\qfigdef{p1_phy_s2_para8_10_4}{Score Advanced - Test 03 (27 Jun 2021)-img/p1_phy_s2_para8_10_4.png}{94.5}{0}
\qfigdef{p1_phy_s2_para11_13_1}{Score Advanced - Test 03 (27 Jun 2021)-img/p1_phy_s2_para11_13_1.png}{83.5}{0}
\qfigdef{p1_phy_s2_para11_13_2}{Score Advanced - Test 03 (27 Jun 2021)-img/p1_phy_s2_para11_13_2.png}{26.3}{0}
\qfigdef{p1_phy_s2_para11_13_3}{Score Advanced - Test 03 (27 Jun 2021)-img/p1_phy_s2_para11_13_3.png}{35.7}{0}
\qfigdef{p1_phy_s2_para11_13_4}{Score Advanced - Test 03 (27 Jun 2021)-img/p1_phy_s2_para11_13_4.png}{9.7}{0}
\qfigdef{p1_phy_s2_para11_13_5}{Score Advanced - Test 03 (27 Jun 2021)-img/p1_phy_s2_para11_13_5.png}{18.8}{0}
\qfigdef{p1_phy_s3_q01_1}{Score Advanced - Test 03 (27 Jun 2021)-img/p1_phy_s3_q01_1.png}{100.7}{0}
\qfigdef{p1_phy_s3_q05_1}{Score Advanced - Test 03 (27 Jun 2021)-img/p1_phy_s3_q05_1.png}{104.3}{0}
\qfigdef{p1_phy_s3_q05_2}{Score Advanced - Test 03 (27 Jun 2021)-img/p1_phy_s3_q05_2.png}{154.3}{0}
\qfigdef{p1_phy_s3_q05_3}{Score Advanced - Test 03 (27 Jun 2021)-img/p1_phy_s3_q05_3.png}{144.3}{0}
\qfigdef{p1_chem_s1_q01_1}{Score Advanced - Test 03 (27 Jun 2021)-img/p1_chem_s1_q01_1.png}{99.2}{0}
\qfigdef{p1_chem_s1_q05_1}{Score Advanced - Test 03 (27 Jun 2021)-img/p1_chem_s1_q05_1.png}{272.0}{1}
\qfigdef{p1_chem_s1_q06_2}{Score Advanced - Test 03 (27 Jun 2021)-img/p1_chem_s1_q06_2.png}{116.7}{0}
\qfigdef{p1_chem_s1_q06_1}{Score Advanced - Test 03 (27 Jun 2021)-img/p1_chem_s1_q06_1.png}{122.6}{0}
\qfigdef{p1_chem_s1_q06_3}{Score Advanced - Test 03 (27 Jun 2021)-img/p1_chem_s1_q06_3.png}{118.2}{0}
\qfigdef{p1_chem_s1_q06_4}{Score Advanced - Test 03 (27 Jun 2021)-img/p1_chem_s1_q06_4.png}{100.0}{0}
\qfigdef{p2_phy_s1_q02_1}{Score Advanced - Test 03 (27 Jun 2021)-img/p2_phy_s1_q02_1.png}{105.1}{0}
\qfigdef{p2_phy_s1_q05_1}{Score Advanced - Test 03 (27 Jun 2021)-img/p2_phy_s1_q05_1.png}{106.0}{0}
\qfigdef{p2_phy_s1_q10_1}{Score Advanced - Test 03 (27 Jun 2021)-img/p2_phy_s1_q10_1.png}{129.9}{0}
\qfigdef{p2_phy_s1_q11_1}{Score Advanced - Test 03 (27 Jun 2021)-img/p2_phy_s1_q11_1.png}{89.5}{0}
\qfigdef{p2_phy_s1_q12_1}{Score Advanced - Test 03 (27 Jun 2021)-img/p2_phy_s1_q12_1.png}{68.0}{0}
\qfigdef{p2_phy_s1_q13_1}{Score Advanced - Test 03 (27 Jun 2021)-img/p2_phy_s1_q13_1.png}{104.0}{0}
\qfigdef{p2_phy_s1_q14_1}{Score Advanced - Test 03 (27 Jun 2021)-img/p2_phy_s1_q14_1.png}{51.3}{0}
\qfigdef{p2_phy_s1_q15_1}{Score Advanced - Test 03 (27 Jun 2021)-img/p2_phy_s1_q15_1.png}{103.8}{0}
\qfigdef{p2_phy_s1_q16_1}{Score Advanced - Test 03 (27 Jun 2021)-img/p2_phy_s1_q16_1.png}{149.0}{0}
\qfigdef{p2_phy_s1_q17_1}{Score Advanced - Test 03 (27 Jun 2021)-img/p2_phy_s1_q17_1.png}{129.5}{0}
\qfigdef{p2_phy_s1_q18_1}{Score Advanced - Test 03 (27 Jun 2021)-img/p2_phy_s1_q18_1.png}{109.5}{0}
\qfigdef{p2_phy_s1_q19_1}{Score Advanced - Test 03 (27 Jun 2021)-img/p2_phy_s1_q19_1.png}{88.6}{0}
\qfigdef{p2_phy_s2_q03_1}{Score Advanced - Test 03 (27 Jun 2021)-img/p2_phy_s2_q03_1.png}{79.0}{0}
\qfigdef{p2_phy_s2_q06_1}{Score Advanced - Test 03 (27 Jun 2021)-img/p2_phy_s2_q06_1.png}{104.0}{0}
\qfigdef{p2_phy_s2_q09_1}{Score Advanced - Test 03 (27 Jun 2021)-img/p2_phy_s2_q09_1.png}{102.2}{0}
\qfigdef{p2_chem_s1_q04_1}{Score Advanced - Test 03 (27 Jun 2021)-img/p2_chem_s1_q04_1.png}{140.0}{0}
\qfigdef{p2_chem_s1_q07_1}{Score Advanced - Test 03 (27 Jun 2021)-img/p2_chem_s1_q07_1.png}{146.7}{0}
\qfigdef{p2_chem_s1_q07_2}{Score Advanced - Test 03 (27 Jun 2021)-img/p2_chem_s1_q07_2.png}{104.0}{0}
\qfigdef{p2_chem_s1_q07_3}{Score Advanced - Test 03 (27 Jun 2021)-img/p2_chem_s1_q07_3.png}{102.4}{0}
\qfigdef{p2_chem_s1_q07_4}{Score Advanced - Test 03 (27 Jun 2021)-img/p2_chem_s1_q07_4.png}{104.0}{0}
\qfigdef{p2_chem_s1_q07_5}{Score Advanced - Test 03 (27 Jun 2021)-img/p2_chem_s1_q07_5.png}{102.5}{0}
\qfigdef{p2_chem_s1_q08_1}{Score Advanced - Test 03 (27 Jun 2021)-img/p2_chem_s1_q08_1.png}{74.9}{0}
\qfigdef{p2_chem_s1_q11_1}{Score Advanced - Test 03 (27 Jun 2021)-img/p2_chem_s1_q11_1.png}{70.5}{0}
\qfigdef{p2_chem_s1_q11_2}{Score Advanced - Test 03 (27 Jun 2021)-img/p2_chem_s1_q11_2.png}{37.4}{0}
\qfigdef{p2_chem_s1_q16_1}{Score Advanced - Test 03 (27 Jun 2021)-img/p2_chem_s1_q16_1.png}{121.7}{0}
\qfigdef{p2_chem_s1_q16_2}{Score Advanced - Test 03 (27 Jun 2021)-img/p2_chem_s1_q16_2.png}{119.3}{0}
\qfigdef{p2_chem_s1_q16_3}{Score Advanced - Test 03 (27 Jun 2021)-img/p2_chem_s1_q16_3.png}{126.8}{0}
\qfigdef{p2_chem_s1_q19_1}{Score Advanced - Test 03 (27 Jun 2021)-img/p2_chem_s1_q19_1.png}{101.2}{0}
\qfigdef{p2_chem_s1_q19_2}{Score Advanced - Test 03 (27 Jun 2021)-img/p2_chem_s1_q19_2.png}{103.3}{0}
\qfigdef{p2_chem_s1_q19_3}{Score Advanced - Test 03 (27 Jun 2021)-img/p2_chem_s1_q19_3.png}{100.7}{0}
\qfigdef{p2_chem_s1_q19_4}{Score Advanced - Test 03 (27 Jun 2021)-img/p2_chem_s1_q19_4.png}{79.0}{0}
\qfigdef{p2_chem_s1_q20_1}{Score Advanced - Test 03 (27 Jun 2021)-img/p2_chem_s1_q20_1.png}{77.6}{0}
\qfigdef{p2_chem_s2_q08_1}{Score Advanced - Test 03 (27 Jun 2021)-img/p2_chem_s2_q08_1.png}{119.9}{0}
\raggedcolumns
\makeatletter\@beginparpenalty=10000\makeatother
\newcolumntype{L}[1]{>{\raggedright\arraybackslash}p{#1}}
\begin{document}
\begin{center}{\Large\bfseries Score Advanced --- Test 03 (27 Jun 2021)}\end{center}
\begin{multicols}{2}
\raggedcolumns
\renewcommand{\paraq}[2]{\par\smallskip\noindent\textbf{Paragraph for Questions #1 - #2}\par}
\papersec{1}
\partsec{1}{PHYSICS}
\sect{SEC-1}
\qn{1} A boat is traveling due east at 12 m/s. A flag on the boat flaps at 53$^\circ$ N of W. Another flag on the shore flaps due north.
\begin{qoptions}
\opt{A}{Speed of wind with respect to ground is 16 m/s}
\opt{B}{Speed of wind with respect to ground is 20 m/s}
\opt{C}{Speed of wind with respect to boat is 20 m/s}
\opt{D}{Speed of wind with respect to boat is 16 m/s}
\end{qoptions}
\qn{2} A block of mass 1 kg is stationary with respect to a conveyor belt that is accelerating with 1 m/s$^2$ upwards at an angle of 30$^\circ$ as shown in figure. What can be the coefficient of static friction between the block and the belt?\par
\qfig{p1_phy_s1_q02_1}
\begin{qoptions}
\opt{A}{0.4}
\opt{B}{0.6}
\opt{C}{0.8}
\opt{D}{1.0}
\end{qoptions}
\qn{3} A ball of mass m collides elastically with an identical ball at rest with some impact parameter.
\begin{qoptions}
\opt{A}{100 \% energy transfer can never take place}
\opt{B}{100 \% energy transfer may take place}
\opt{C}{angle of divergence between the two balls must be 90$^\circ$}
\opt{D}{angle of divergence between the two balls depend on impact parameter}
\end{qoptions}
\qn{4} A billiard ball initially at rest is given a sharp blow by a cue stick. The impulse is horizontal and is applied at a distance 2R/3 below the centreline, as shown in figure. The initial speed of the ball is $\upsilon_0$, and the coefficient of kinetic friction is $\mu_k$.\par
\qfig{p1_phy_s1_q04_1}
\begin{qoptions}
\opt{A}{Initially kinetic friction acts in $-\hat{i}$ direction}
\opt{B}{Initially kinetic friction in $\hat{i}$ direction.}
\opt{C}{Ball instantaneously starts pure rolling.}
\opt{D}{Initial angular velocity of ball is $\dfrac{5v_0}{3R}$}
\end{qoptions}
\qn{5} Which of the following statement(s) is/are true ?
\begin{qoptions}
\opt{A}{Maximum velocity of photoelectrons depends on frequency and intensity of the incident light}
\opt{B}{Maximum velocity of photoelectrons depends only on the frequency of the incident light}
\opt{C}{Photoelectric current increases with increase in intensity of incident light by increasing only the frequency.}
\opt{D}{Photoelectric current increases with increase in intensity of incident light by increasing the number of photons.}
\end{qoptions}
\qn{6} Figure shows a small particle in a box which bounces back and forth at constant speed. The reflections of particle from the wall creates a standing wave analogous to standing wave on a string tied at both ends. Since a standing wave confined to a region can have only selected wavelengths, momentum of the particle is quantized. We can safely assume that such a particle only has kinetic energy. This energy must also be quantized.\par
\qfig{p1_phy_s1_q06_1}
\begin{qoptions}
\opt{A}{The momentum of particle in n$^{th}$ mode of standing wave is $\dfrac{nh}{2L}$}
\opt{B}{Particle's energy in n$^{th}$ state is $\dfrac{h^2 n^2}{8mL^2}$}
\opt{C}{Matter wave in this case will be electromagnetic in nature.}
\opt{D}{None of these}
\end{qoptions}
\qn{7} The $^{238}$U nucleus has a binding energy of about 7.6 MeV per nucleon. If the nucleus were to fission into two equal fragments, each would have a kinetic energy of just over 100 MeV. From this, it can be concluded that
\begin{qoptions}
\opt{A}{nuclei near A = 119 has mass less than half that of $^{238}$U.}
\opt{B}{nuclei near A = 119 have masses greater than half that of $^{238}$U}
\opt{C}{nuclei near A = 119 must be bound by about 6.7 MeV/ nucleon}
\opt{D}{nuclei near A = 119 must be bound by about 8.5 MeV/ nucleon}
\end{qoptions}
\sect{SEC-2}
\paraq{8}{10}
\noindent{\footnotesize\setlength{\tabcolsep}{2pt}\begin{tabular}{|L{.8cm}|L{3.8cm}|L{1.2cm}|L{1.2cm}|}\hline
\multicolumn{2}{|L{4.6cm}|}{Column--1} & Column--2 & Column--3\\\hline
\multicolumn{2}{|L{4.6cm}|}{Mass of the Object is `m' situation is given at t = 0} & Kinetic energy at the time when slipping ceases. & Time till slipping ceases.\\\hline
(I) & \qfigi{p1_phy_s2_para8_10_1} & (i) Zero & (P) $\dfrac{v_0}{\mu g}$\\\hline
(II) & \qfigi{p1_phy_s2_para8_10_2} & (ii) $\dfrac{1}{12}mv_0^2$ & (Q) $\dfrac{2v_0}{3\mu g}$\\\hline
(III) & \qfigi{p1_phy_s2_para8_10_3} & (iii) $\dfrac{9}{70}mv_0^2$ & (R) $\dfrac{4v_0}{7\mu g}$\\\hline
(IV) & \qfigi{p1_phy_s2_para8_10_4} & (iv) $\dfrac{mv_0^2}{30}$ & (S) $\dfrac{4v_0}{5\mu g}$\\\hline
\end{tabular}}\par
\qn{1} Name the object that has maximum loss in kinetic energy, its kinetic energy after slipping ceases and the time till slipping ceases.
\begin{qoptions}
\opt{A}{(II) (iii) (Q)}
\opt{B}{(I) (i) (P)}
\opt{C}{(III) (iii) (R)}
\opt{D}{(IV) (iv) (S)}
\end{qoptions}
\qn{2} Name the object has minimum loss in kinetic energy during slipping, its kinetic energy after slipping ceases and the time till slipping ceases :-
\begin{qoptions}
\opt{A}{(II) (ii) (P)}
\opt{B}{(I) (i) (S)}
\opt{C}{(III) (iii) (R)}
\opt{D}{(IV) (iv) (Q)}
\end{qoptions}
\qn{3} Name the object, which stops when slipping ceases, its kinetic energy at the time of slipping ceases and time till slipping ceases :-
\begin{qoptions}
\opt{A}{(II) (ii) (R)}
\opt{B}{(III) (iii) (S)}
\opt{C}{(IV) (iv) (Q)}
\opt{D}{(I) (i) (P)}
\end{qoptions}
\paraq{11}{13}A particle is going around a circle with constant speed v of radius r. At t = 0 particle is at A. Particle takes 4 seconds to complete one round.\par
\qfig{p1_phy_s2_para11_13_1}
\noindent{\footnotesize\setlength{\tabcolsep}{2pt}\begin{tabular}{|L{.8cm}|L{1.9cm}|L{1.3cm}|L{2.2cm}|}\hline
\multicolumn{2}{|L{2.7cm}|}{Column--1} & Column--2 & Column--3\\\hline
\multicolumn{2}{|L{2.7cm}|}{magnitude of displacement (in m)} & Direction of displacement vector & Magnitude of change in velocity (in m/s)\\\hline
(I) & 2R & (i) \qfigi{p1_phy_s2_para11_13_2} & (P) $\sqrt{2}$V\\\hline
(II) & $2\sqrt{2}$R & (ii) \qfigi{p1_phy_s2_para11_13_3} & (Q) 2V\\\hline
(III) & $\sqrt{2}$R & (iii) \qfigi{p1_phy_s2_para11_13_4} & (R) $2\sqrt{2}$V\\\hline
(IV) & R & (iv) \qfigi{p1_phy_s2_para11_13_5} & (S) V\\\hline
\end{tabular}}\par
\qn{4} Between A \& B :-
\begin{qoptions}
\opt{A}{(II) (iv) (P)}
\opt{B}{(I) (iii) (S)}
\opt{C}{(III) (iv) (P)}
\opt{D}{(III) (i) (P)}
\end{qoptions}
\qn{5} Between B \& D :-
\begin{qoptions}
\opt{A}{(I) (iii) (R)}
\opt{B}{(I) (ii) (Q)}
\opt{C}{(II) (ii) (R)}
\opt{D}{(I) (ii) (P)}
\end{qoptions}
\qn{6} From t = 2 to t = 4 sec :-
\begin{qoptions}
\opt{A}{(I) (iii) (Q)}
\opt{B}{(I) (ii) (R)}
\opt{C}{(III) (iii) (Q)}
\opt{D}{(IV) (i) (P)}
\end{qoptions}
\sect{SEC-3}
\qn{1} Two blocks of masses 3 kg and 6 kg are connected by an ideal relaxed spring and are placed on a frictionless horizontal surface. The 3 kg block is imparted a speed of 2 m/s towards left. (k = 2 N/m)\par
\qfig{p1_phy_s3_q01_1}
\noindent{\footnotesize\setlength{\tabcolsep}{2pt}\begin{tabular}{|L{.7cm}|L{4.0cm}|L{.7cm}|L{1.3cm}|}\hline
 & List-I & & List-II\\\hline
(P) & The speed of 3 kg block is $\dfrac{2}{3}$ m/s at time & (1) & $\pi$ sec\\\hline
(Q) & The speed of 6 kg block is $\dfrac{4}{3}$ m/s at time & (2) & $\dfrac{\pi}{3}$ sec\\\hline
(R) & The speed of 3 kg block is $\dfrac{4}{3}$ m/s at time & (3) & $\dfrac{\pi}{4}$ sec\\\hline
\end{tabular}}\par
\qn{2} Three physical quantities are listed in list--I and their values are listed in list--II in random order.\par
\noindent{\footnotesize\setlength{\tabcolsep}{2pt}\begin{tabular}{|L{.7cm}|L{3.9cm}|L{.7cm}|L{1.4cm}|}\hline
 & List-I & & List-II\\\hline
(P) & Binding energy of heavy nuclei per nucleon & (1) & 10 keV\\\hline
(Q) & X-ray photon energy & (2) & 7 MeV\\\hline
(R) & Photon energy of visible light & (3) & 2 eV\\\hline
\end{tabular}}\par
\qn{3} Match the statements given in List-I with their corresponding possible results in List-II.\par
\noindent{\footnotesize\setlength{\tabcolsep}{2pt}\begin{tabular}{|L{.7cm}|L{4.0cm}|L{.7cm}|L{1.3cm}|}\hline
 & List-I & & List-II\\\hline
(P) & If photons of ultraviolet light of energy 12 eV are incident on a metal surface of work function of 3eV, then the stopping potential (in eV) will be & (1) & 8\\\hline
(Q) & The ratio of wavelengths of K$_\alpha$ lines of two elements is $\left(\dfrac{85}{81}\right)^2$. Number of elements having atomic numbers between these elements will be & (2) & 3\\\hline
(R) & If 20 gm of a radioactive substance due to radioactive decay reduces to 10gm in 4 minutes, then in what time (in minutes) 80 gm of the same substance will reduce to 20 gm. & (3) & 9\\\hline
\end{tabular}}\par
\qn{4} Two particles A and B moving in x-y plane are at origin at t = 0 sec. The initial velocity vectors of A and B are $\vec{u}_A = 8\hat{i}$ m/s and $\vec{u}_B = 8\hat{j}$ m/s. The acceleration of A and B are constant and are $\vec{a}_A = -2\hat{i}$ m/s$^2$ and $\vec{a}_B = -2\hat{j}$ m/s$^2$. List-I gives certain statements regarding particle A and B. List-II gives corresponding results. Match the statements in List-I with corresponding results in List-II.\par
\noindent{\footnotesize\setlength{\tabcolsep}{2pt}\begin{tabular}{|L{.7cm}|L{4.0cm}|L{.7cm}|L{1.3cm}|}\hline
 & List-I & & List-II\\\hline
(P) & The time (in seconds) at which velocity of A relative to B is zero & (1) & 16$\sqrt{2}$\\\hline
(Q) & The distance (in metres) between A and B when their relative velocity is zero & (2) & 8$\sqrt{2}$\\\hline
(R) & The magnitude of relative velocity of A and B at the instant they are at same position. & (3) & 4\\\hline
\end{tabular}}\par
\qn{5} Match the value of moment of inertia of structures in List-I with corresponding results in List-II.\par
\noindent{\footnotesize\setlength{\tabcolsep}{2pt}\begin{tabular}{|L{.6cm}|L{4.3cm}|L{.6cm}|L{1.5cm}|}\hline
 & List-I & & List-II\\\hline
(P) & \qfigi{p1_phy_s3_q05_1} & (1) & $\dfrac{8MR^2}{11}$\\\hline
(Q) & \qfigi{p1_phy_s3_q05_2} & (2) & $\dfrac{MR^2}{12}$\\\hline
(R) & \qfigi{p1_phy_s3_q05_3} & (3) & $\dfrac{13MR^2}{8}$\\\hline
\end{tabular}}\par
\partsec{2}{CHEMISTRY}
\sect{SEC-1}
\qn{1} Two gases A \& B are mixed under constant temperature condition. Initial volume of gas A is 2 litre \& that of gas B is 5 litre. When stop cock is opened then which of the following statements is/are correct (initials moles of gas A = 2 \& gas B = 5)-\par
\qfig{p1_chem_s1_q01_1}
\begin{qoptions}
\opt{A}{$\Delta S_{mix} = \ln\left[\left(\dfrac{7}{2}\right)^{2R} \times \left(\dfrac{7}{5}\right)^{5R}\right]$}
\opt{B}{$\Delta S_{mix}$ = 0, if gas `A' \& `B' are same and equal volumes of them are mixed}
\opt{C}{$\Delta S_{mix}$ = (n$_1$ + n$_2$) R ln2 if equal volumes of same gases are mixed}
\opt{D}{$\Delta S_{mix}$ is maximum when equal moles of two gases are mixed. (Gas `A' \& `B' are different)}
\end{qoptions}
\qn{2} Which of the following statement is/are correct?
\begin{qoptions}
\opt{A}{Sorel cement contains MgO \& MgCl$_2$}
\opt{B}{Silicates of calcium are important ingredients of portland cement}
\opt{C}{Be \& Mg dissolve in liquid ammonia to give deep blue black solution}
\opt{D}{K$_2$CO$_3$ and Na$_2$CO$_3$ is used as fusion mixture in laboratory}
\end{qoptions}
\qn{3} Which of the following statements are correct?
\begin{qoptions}
\opt{A}{K$_2$Cr$_2$O$_7$ is deliquescent but Na$_2$Cr$_2$O$_7$ is not a deliquescent}
\opt{B}{In CrO$_4^{2-}$ \& Cr$_2$O$_7^{2-}$ geometry around Cr is same}
\opt{C}{CrO$_4^{2-}$ \& Cr$_2$O$_7^{2-}$ produce orange \& yellow colour respectively due to charge transfer}
\opt{D}{In Deacon process CuCl$_2$ is used as a catalyst in making Cl$_2$ from HCl}
\end{qoptions}
\qn{4} Which of the following compound may be used to remove temporary hardness?
\begin{qoptions}
\opt{A}{Ca(OH)$_2$}
\opt{B}{Na$_6$P$_6$O$_{18}$}
\opt{C}{Na$_2$Al$_2$Si$_2$O$_8$.xH$_2$O}
\opt{D}{Na$_2$CO$_3$.10H$_2$O}
\end{qoptions}
\qfullbegin
\qn{5} Consider the following reaction :-\par
\qfig{p1_chem_s1_q05_1}
\qfullsplit
Which of the following statement is/are correct?
\begin{qoptions}
\opt{A}{`Z' is isomorphous with Na$_2$SO$_4$.10H$_2$O}
\opt{B}{In basic medium MnO$_4^{\ominus}$ produce `X' in the form of brown ppt. when it is act as oxidising agent}
\opt{C}{`Y' produce `X' in the form of brown ppt. \& KMnO$_4$ with water}
\opt{D}{`Y' produce `X' in the form of brown ppt. \& KMnO$_4$ with water \& CO$_2$}
\end{qoptions}
\qfullend
\qn{6} Which one of the following reaction will produce hydrocarbon as final product?
\begin{qoptions}
\opt{A}{\qfigi{p1_chem_s1_q06_2}}
\opt{B}{\qfigi{p1_chem_s1_q06_1}}
\opt{C}{\qfigi{p1_chem_s1_q06_3}}
\opt{D}{\qfigi{p1_chem_s1_q06_4}}
\end{qoptions}
\qn{7} Which Amino acid(s) is/are essential amino acid -
\begin{qoptions}
\opt{A}{Glycine}
\opt{B}{Phenyl alanine}
\opt{C}{Valine}
\opt{D}{Threonine}
\end{qoptions}
\sect{SEC-2}
\paraq{8}{10}
\noindent{\footnotesize\setlength{\tabcolsep}{2pt}\begin{tabular}{|L{1.9cm}|L{2.2cm}|L{2.7cm}|}\hline
Column-I & Column-II & Column-III\\\hline
(A) V$_m$ = 50 L & (P) $\dfrac{P}{RT} = \dfrac{1}{75}$ & (I) Z $<$ 1\\\hline
(B) V$_m$ = 80 L & (Q) $\dfrac{P}{RT} = \dfrac{1}{100}$ & (II) Z $>$ 1\\\hline
(C) V$_m$ = 120 L & (R) $\dfrac{P}{RT} = \dfrac{1}{40}$ & (III) Attractive forces dominant\\\hline
(D) V$_m$ = 20 L & (S) $\dfrac{P}{RT} = \dfrac{1}{140}$ & (IV) Molecular size dominant\\\hline
\end{tabular}}\par
\qn{1} If rate of diffusion of a real gas is 2/5 times of rate of diffusion of He and density of real gas = 0.5 gm/l then which is incorrect matching
\begin{qoptions}
\opt{A}{A, P, I}
\opt{B}{A, Q, I}
\opt{C}{A, R, I}
\opt{D}{A, P, III}
\end{qoptions}
\qn{2} If molar mass of a real gas is 450 \& density is 3.75 gm/litre then which is correct matching
\begin{qoptions}
\opt{A}{C, P, I}
\opt{B}{C, Q, I}
\opt{C}{C, S, IV}
\opt{D}{C, R, II}
\end{qoptions}
\qn{3} If molar mass of a real gas is 500 \& 10 litre it has a weight of 62.5 gms then which INCORRECT matching
\begin{qoptions}
\opt{A}{B, P, II}
\opt{B}{B, R, IV}
\opt{C}{B, S, I}
\opt{D}{B, R, III}
\end{qoptions}
\paraq{11}{13}
\noindent{\scriptsize\setlength{\tabcolsep}{2pt}\begin{tabular}{|L{.6cm}|L{1.4cm}|L{.6cm}|L{2.1cm}|L{.6cm}|L{1.6cm}|}\hline
 & Column-I (Carbohydrate) & & Column-II (Monomer unit) & & Column-III (Nnature of carbohydrate)\\\hline
(P) & Cellulose & (i) & $\alpha$-Glucose & (1) & Reducing sugar\\\hline
(Q) & Sucrose & (ii) & $\beta$-Glucose & (2) & Non reducing\\\hline
(R) & Starch & (iii) & Glucopyranose \& Fructofuranose & (3) & Mono-saccharides\\\hline
(S) & Lactose & (iv) & Glucose \& $\beta$-Galactose & (4) & Poly-saccharides\\\hline
\end{tabular}}\par
\qn{4} Which of the following combination is correct for Cellulose ?
\begin{qoptions}
\opt{A}{ii , 2 , 4}
\opt{B}{i , 2 , 4}
\opt{C}{iv , 4}
\opt{D}{ii , 2 , 3}
\end{qoptions}
\qn{5} Which of the following is not correctly match with their monomer ?
\begin{qoptions}
\opt{A}{P , ii}
\opt{B}{Q , iii}
\opt{C}{R , i}
\opt{D}{S , iii}
\end{qoptions}
\qn{6} Which of the following is non reducing sugar ?
\begin{qoptions}
\opt{A}{$\alpha$-Glucose}
\opt{B}{$\beta$-Glucose}
\opt{C}{Galactose}
\opt{D}{Starch}
\end{qoptions}
\sect{SEC-3}
\qn{1} 1 mole ideal gas undergoes change in state from 2 atm, 5 litre, 300K) to (P$_f$ , V$_f$ \& T$_f$)\par
\noindent{\footnotesize\setlength{\tabcolsep}{2pt}\begin{tabular}{|L{.7cm}|L{2.9cm}|L{.7cm}|L{2.4cm}|}\hline
(P) & If w = $-$6.93 atm - litre & (1) & Process is irrreversible isothermal at P$_{ext}$= 1atm\\\hline
(Q) & If w = $-$ 5 atm - litre & (2) & Process is isobaric at 2 atm pressure with V$_f$ = 20 litre\\\hline
(R) & If w = $-$ 30 atm - litre & (3) & Process is reversible isothermal with P$_f$ = 1 atm\\\hline
\end{tabular}}\par
\qn{2} \noindent{\footnotesize\setlength{\tabcolsep}{2pt}\begin{tabular}{|L{1.7cm}|L{.75cm}|L{.75cm}|L{.75cm}|L{.75cm}|L{.75cm}|L{.75cm}|}\hline
Gases & A(g) & B(g) & C(g) & D(g) & E(g) & F(g)\\\hline
$\Delta H^0_{f\,kJ/mole}$ & 40 & 60 & 30 & 20 & 80 & 100\\\hline
\end{tabular}}\par
\noindent{\footnotesize\setlength{\tabcolsep}{2pt}\begin{tabular}{|L{.7cm}|L{3.0cm}|L{.7cm}|L{2.3cm}|}\hline
(P) & A(g) $\to$ B(g) + C(g) & (1) & Spontaneous at all temperature\\\hline
(Q) & C(g) + E(g) $\to$ D(g) & (2) & Spontaneous at low temperature\\\hline
(R) & F(g) $\to$ B(g) + C(g) & (3) & Spontaneous at high temperature\\\hline
\end{tabular}}\par
\qn{3} Match given properties with elements\par
\noindent{\footnotesize\setlength{\tabcolsep}{2pt}\begin{tabular}{|L{.6cm}|L{3.7cm}|L{.6cm}|L{1.7cm}|}\hline
 & List-I (Properties ) & & List-I (Elements)\\\hline
P & Produce stable M$_2$O type oxide with limited supply of air & 1 & Li, Be, Al\\\hline
Q & Chloride of elements produce red colour in oxidising flame & 2 & Li, Ca, Sr\\\hline
R & Chlorides of such elements having significant covalent character or covalent nature in anhydrous form & 3 & Na, Li, Cs\\\hline
\end{tabular}}\par
\qn{4} Match the colour with given compounds\par
\noindent{\footnotesize\setlength{\tabcolsep}{2pt}\begin{tabular}{|L{.6cm}|L{2.8cm}|L{.6cm}|L{2.4cm}|}\hline
 & List-I (Compounds) & & List-II (Colour)\\\hline
P & Mn$_2$O$_7$ & 1 & Orange\\\hline
Q & (NH$_4$)$_2$Cr$_2$O$_7$ & 2 & Green\\\hline
R & CrO$_2$Cl$_2$ & 3 & Reddish Brown\\\hline
\end{tabular}}\par
\qn{5} \noindent{\footnotesize\setlength{\tabcolsep}{2pt}\begin{tabular}{|L{.6cm}|L{3.2cm}|L{.6cm}|L{2.4cm}|}\hline
 & List-I & & List-II\\\hline
P & Synthetic condensation copolymer & 1 & Buna-N\\\hline
Q & Distinguish between aldehyde \& ketones & 2 & Terylene\\\hline
R & Addition copolymer & 3 & Tollen's reagent\\\hline
\end{tabular}}\par
\partsec{3}{MATHEMATICS}
\sect{SEC-1}
\qn{1} Let M = $\begin{pmatrix} 1 & 3\\ 0 & 2 \end{pmatrix}$ and I + C$_1$M + C$_2$M$^2$ = $\begin{pmatrix} 0 & 0\\ 0 & 0 \end{pmatrix}$, then answer the following (C$_1$, C$_2$ $\in$ R and I is an identity matrix of order `2')
\begin{qoptions}
\opt{A}{C$_1$ = $\dfrac{1}{2}$}
\opt{B}{C$_1$ = $-\dfrac{3}{2}$}
\opt{C}{C$_2$ = $\dfrac{1}{2}$}
\opt{D}{C$_2$ = $-\dfrac{3}{2}$}
\end{qoptions}
\qn{2} Let A = $\begin{vmatrix} a & a + d & a + 3d\\ a + d & a + 2d & a\\ a + 2d & a & a + d \end{vmatrix}$, then
\begin{qoptions}
\opt{A}{A = 0}
\opt{B}{A depends on a}
\opt{C}{A depends on d}
\opt{D}{A is independent of a and d}
\end{qoptions}
\qn{3} Area bounded by the curve 2$|$x$|$ + 3$|$y$|$ $\leq$ 6 is n, then
\begin{qoptions}
\opt{A}{n $\geq$ 12}
\opt{B}{n $\leq$ 10}
\opt{C}{n $>$ 12}
\opt{D}{n $>$ 10}
\end{qoptions}
\qn{4} Let $f(x) = \begin{cases} x^{2\alpha + 1}\ell nx & x > 0\\ 0 & x = 0 \end{cases}$. If $f(x)$ satisfies rolle's theorem in [0,1], then $\alpha$ can be-
\begin{qoptions}
\opt{A}{$-\dfrac{1}{2}$}
\opt{B}{$-\dfrac{1}{3}$}
\opt{C}{$-\dfrac{1}{4}$}
\opt{D}{$-1$}
\end{qoptions}
\qn{5} Let A = $\begin{pmatrix} 2 & 2 & 1\\ 2 & 5 & 2\\ 1 & 2 & 2 \end{pmatrix}$ and B = $\begin{pmatrix} -x & -y & z\\ 0 & y & 2z\\ x & -y & z \end{pmatrix}$ where x,y,z $\in$ R. If B$^T$AB = $\begin{pmatrix} 8 & 0 & 0\\ 0 & 27 & 0\\ 0 & 0 & 42 \end{pmatrix}$, then xyz can be
\begin{qoptions}
\opt{A}{4}
\opt{B}{$-$4}
\opt{C}{6}
\opt{D}{$-$6}
\end{qoptions}
\qn{6} Consider two curves y = e$^x$, y = e$^{a-x}$, where a is positive number. Let S be the area of the part surrounded by C$_1$,C$_2$ and the y-axis, if $\displaystyle\lim_{a \to 0}\dfrac{S}{a^2}$ is $\dfrac{1}{n}$, then
\begin{qoptions}
\opt{A}{n is the first even composite number}
\opt{B}{number of positive divisors of n is 3}
\opt{C}{n is a perfect square}
\opt{D}{n is an odd number}
\end{qoptions}
\qn{7} The straight line which is both tangent and normal to the curve x = 3t$^2$, y = 2t$^3$ is
\begin{qoptions}
\opt{A}{y + $\sqrt{3}$ (x $-$ 1) = 0}
\opt{B}{y $-$ $\sqrt{3}$ (x $-$ 1) = 0}
\opt{C}{y + $\sqrt{2}$ (x $-$ 2) = 0}
\opt{D}{y $-$ $\sqrt{2}$ (x $-$ 2) = 0}
\end{qoptions}
\sect{SEC-2}
\paraq{8}{10}Column-I represents 2 curves.\par
Column-II Angle of intersection of between given curves.\par
Column-III Number of points of intersection of given curves.\par
\noindent{\footnotesize\setlength{\tabcolsep}{2pt}\begin{tabular}{|L{.7cm}|L{2.6cm}|L{.8cm}|L{1.5cm}|L{.6cm}|L{.6cm}|}\hline
(I) & y = sinx and y = cosx both in (0,2$\pi$) & (i) & tan$^{-1}$2$\sqrt{2}$ & (P) & 2\\\hline
(II) & x$^2$ = 4y, y = $\dfrac{8}{x^2 + 4}$ & (ii) & $\dfrac{\pi}{2}$ & (Q) & 4\\\hline
(III) & $\dfrac{x^2}{18} + \dfrac{y^2}{8} = 1$, y = 0 & (iii) & tan$^{-1}$3 & (R) & 0\\\hline
(IV) & xy = 1, x$^2$ $-$ y$^2$ = 5 & (iv) & tan$^{-1}$5 & (S) & 8\\\hline
\end{tabular}}\par
\qn{1} Which of the following options is the only CORRECT combination ?
\begin{qoptions}
\opt{A}{(I) (i) (P)}
\opt{B}{(II) (ii) (P)}
\opt{C}{(III) (iii) (P)}
\opt{D}{(IV) (iv) (P)}
\end{qoptions}
\qn{2} Which of the following options is the only CORRECT combination ?
\begin{qoptions}
\opt{A}{(I) (ii) (P)}
\opt{B}{(II) (iv) (P)}
\opt{C}{(III) (ii) (P)}
\opt{D}{(IV) (iv) (Q)}
\end{qoptions}
\qn{3} Which of the following options is the only CORRECT combination ?
\begin{qoptions}
\opt{A}{(I) (ii) (P)}
\opt{B}{(I) (ii) (Q)}
\opt{C}{(I) (iii) (Q)}
\opt{D}{(IV) (ii) (P)}
\end{qoptions}
\paraq{11}{13}Column-I : y = $f(x)$ with respective domain.\par
Column-II : Interval in which $f(x)$ is increasing.\par
Column-III : Interval in which $f(x)$ is decreasing.\par
\noindent{\scriptsize\setlength{\tabcolsep}{2pt}\begin{tabular}{|L{.6cm}|L{2.0cm}|L{.7cm}|L{1.5cm}|L{.6cm}|L{1.5cm}|}\hline
(I) & y = sinx + cosx in $\left(0, \dfrac{\pi}{2}\right)$ & (i) & $\left(0, \dfrac{\pi}{4}\right)$ & (P) & $\left(0, \dfrac{\pi}{4}\right)$\\\hline
(II) & y = sin2x in $\left(0, \dfrac{\pi}{2}\right)$ & (ii) & $\left(\dfrac{\pi}{4}, \dfrac{\pi}{2}\right)$ & (Q) & $\left(\dfrac{\pi}{4}, \dfrac{\pi}{2}\right)$\\\hline
(III) & y = sin$^5$x + cos$^5$x $-$ 1 in $\left(0, \dfrac{\pi}{2}\right)$ & (iii) & $\left(0, \dfrac{\pi}{3}\right)$ & (R) & $\left(0, \dfrac{1}{e}\right)$\\\hline
(IV) & y = x$^x$ in $\left(0, \dfrac{\pi}{2}\right)$ & (iv) & $\left(\dfrac{1}{e}, \dfrac{\pi}{2}\right)$ & (S) & $\left(\dfrac{1}{e}, \dfrac{\pi}{2}\right)$\\\hline
\end{tabular}}\par
\qn{4} Which of the following options is the only CORRECT combination ?
\begin{qoptions}
\opt{A}{(I) (i) (S)}
\opt{B}{(II) (iii) (S)}
\opt{C}{(III) (ii) (P)}
\opt{D}{(I) (ii) (P)}
\end{qoptions}
\qn{5} Which of the following options is the only CORRECT combination ?
\begin{qoptions}
\opt{A}{(I) (iii) (P)}
\opt{B}{(II) (ii) (P)}
\opt{C}{(III) (iii) (P)}
\opt{D}{(IV) (iv) (R)}
\end{qoptions}
\qn{6} Which of the following options is the only CORRECT combination ?
\begin{qoptions}
\opt{A}{(I) (iv) (R)}
\opt{B}{(II) (i) (Q)}
\opt{C}{(III) (iii) (Q)}
\opt{D}{(II) (ii) (P)}
\end{qoptions}
\sect{SEC-3}
\qn{1} Let $\displaystyle\lim_{T \to \infty}\dfrac{1}{T}\int_0^T (\sin x + \sin ax)^2\,dx$ = L, then\par
\noindent{\footnotesize\setlength{\tabcolsep}{2pt}\begin{tabular}{|L{.7cm}|L{3.8cm}|L{.7cm}|L{1.3cm}|}\hline
 & List-I & & List-II\\\hline
(P) & for a = 0 the value of L is & (1) & 0\\\hline
(Q) & for a = 1 the value of L is & (2) & $\dfrac{1}{2}$\\\hline
(R) & for a = $-$1 the value of L is & (3) & 2\\\hline
\end{tabular}}\par
\qn{2} \noindent{\footnotesize\setlength{\tabcolsep}{2pt}\begin{tabular}{|L{.7cm}|L{4.2cm}|L{.7cm}|L{1.2cm}|}\hline
 & List-I & & List-II\\\hline
(P) & Value of definite integral $\displaystyle\int_{-1}^{1}\dfrac{dx}{\left(1 + x^3 + \sqrt{1 + x^6}\right)}$ & (1) & 3\\\hline
(Q) & Let $f(x) = \dfrac{x - 1}{x^2}$ ; x = k is point of inflection, then k is & (2) & 1\\\hline
(R) & $\displaystyle\int_1^e \dfrac{dx}{\ell n\left(x^x e^x\right)}$ equals & (3) & $\ell$n2\\\hline
\end{tabular}}\par
\qn{3} \noindent{\footnotesize\setlength{\tabcolsep}{2pt}\begin{tabular}{|L{.7cm}|L{4.0cm}|L{.7cm}|L{1.4cm}|}\hline
 & List-I & & List-II\\\hline
(P) & $\begin{vmatrix} 1 & \cos\alpha & \cos\beta\\ \cos\alpha & 1 & \cos\gamma\\ \cos\beta & \cos\gamma & 1 \end{vmatrix}$ = $\begin{vmatrix} 0 & \cos\alpha & \cos\beta\\ \cos\alpha & 0 & \cos\gamma\\ \cos\beta & \cos\gamma & 0 \end{vmatrix}$, then $\sin^2\alpha + \sin^2\beta + \sin^2\gamma$ = & (1) & 2\\\hline
(Q) & If A = $\begin{pmatrix} 3 & 1\\ -1 & 1 \end{pmatrix}$ and a matrix is defined as C = BAB$^{-1}$(B$^{-1}$A$^T$B), then detC is equal to & (2) & 4\\\hline
(R) & If A = $\begin{pmatrix} 1 & 1\\ -1 & 1 \end{pmatrix}$ and A$^4$ = $-\lambda$I, then $\lambda$ equals to & (3) & 16\\\hline
\end{tabular}}\par
\qn{4} \noindent{\footnotesize\setlength{\tabcolsep}{2pt}\begin{tabular}{|L{.7cm}|L{4.2cm}|L{.7cm}|L{1.2cm}|}\hline
 & List-I & & List-II\\\hline
(P) & If I = $\displaystyle\int_0^1 \ell n\left(\dfrac{1}{x} - 1\right)$, then I equals to & (1) & $\dfrac{1}{3}\ell$n9\\\hline
(Q) & $\displaystyle\int_0^{10\pi}\left(\lim_{x \to y}\left(\dfrac{\sin x - \sin y}{x - y}\right)\right)$ dy is & (2) & 0\\\hline
(R) & {\scriptsize$\displaystyle\lim_{n \to \infty}$ $\left(\frac{1^2}{1^3 + n^3} + \frac{2^2}{2^3 + n^3} + \ldots + \frac{4}{9n}\right)$} equal to & (3) & 1\\\hline
\end{tabular}}\par
\qn{5} \noindent{\footnotesize\setlength{\tabcolsep}{2pt}\begin{tabular}{|L{.7cm}|L{4.0cm}|L{.7cm}|L{1.4cm}|}\hline
 & List-I & & List-II\\\hline
(P) & A curve passing through (2,3) having the property that length of the radius vector of any of point P is equal to the length of tangent drawn at this point, curve can be & (1) & Straight line\\\hline
(Q) & A curve passing through (1,1) having the property that any tangent intersects the y-axis at the point which is equidistant from the point of tangency and the origin curve can be & (2) & Circle\\\hline
(R) & A curve passing through (1,0) for which length of normal is equal to radius vector, curve can be & (3) & Parabola\\\hline
\end{tabular}}\par
\raggedcolumns
\papersec{2}
\partsec{1}{PHYSICS}
\sect{SEC-A}
\qn{1} In certain electronic transition from quantum level n to ground state in atomic hydrogen in one or more steps no line belonging to Brackett series is observed. The wave numbers which may be observed in Balmer series is :-
\begin{qoptions}
\opt{A}{$\dfrac{8R}{9},\ \dfrac{5R}{36}$}
\opt{B}{$\dfrac{3R}{16},\ \dfrac{8R}{9}$}
\opt{C}{$\dfrac{5R}{36},\ \dfrac{3R}{16}$}
\opt{D}{$\dfrac{3R}{4},\ \dfrac{3R}{16}$}
\end{qoptions}
\qn{2} A small ball is given a velocity $\sqrt{\dfrac{7}{2}g\ell}$ at the lowest point of verticle circle. Find the maximum height reached by the ball :-\par
\qfig{p2_phy_s1_q02_1}
\begin{qoptions}
\opt{A}{$\dfrac{4\ell}{5}$}
\opt{B}{$\dfrac{5\ell}{4}$}
\opt{C}{$\dfrac{3\ell}{2}$}
\opt{D}{$\dfrac{27\ell}{16}$}
\end{qoptions}
\qn{3} In Bohr's theory for hydrogen-like atoms
\begin{qoptions}
\opt{A}{For the higher atomic number, radius of the n$^{\text{th}}$ orbit is higher.}
\opt{B}{For the higher quantum number, velocity of electron in that orbit is greater.}
\opt{C}{In some hypothetical condition if mass of electron would have been higher then electron velocity in that orbit would have been same.}
\opt{D}{For the higher quantum number, time period of an electron revolving around nucleus is smaller.}
\end{qoptions}
\qn{4} The K, L and M energy levels of platinum lie roughly at 78, 12 and 3 keV respectively. The ratio of wavelength of K$_\alpha$ line to that of K$_\beta$ line in X-ray spectrum is :
\begin{qoptions}
\opt{A}{$\dfrac{22}{3}$}
\opt{B}{$\dfrac{3}{22}$}
\opt{C}{$\dfrac{22}{25}$}
\opt{D}{$\dfrac{25}{22}$}
\end{qoptions}
\qn{5} The normal reaction force between the blocks and direction of acceleration are :-\par
\qfig{p2_phy_s1_q05_1}
\begin{qoptions}
\opt{A}{8 N, up the incline}
\opt{B}{8N, down the incline}
\opt{C}{5N, up the incline}
\opt{D}{5N, down the incline}
\end{qoptions}
\qn{6} A normal stationary lithium atom releases an alpha particle with velocity $\left(2\hat{i}+3\hat{j}+7\hat{k}\right)$ km/s. What will be the velocity of the daughter element just after the release of the alpha particle?
\begin{qoptions}
\opt{A}{$(1/2)\left(2\hat{i}+3\hat{j}+7\hat{k}\right)$ km/s}
\opt{B}{$(-4/3)\left(2\hat{i}+3\hat{j}+7\hat{k}\right)$ km/s}
\opt{C}{$(-1/4)\left(2\hat{i}+3\hat{j}+7\hat{k}\right)$ km/s}
\opt{D}{$(1/5)\left(2\hat{i}+3\hat{j}+7\hat{k}\right)$ km/s}
\end{qoptions}
\qn{7} The emitted $_{-1}\beta^0$ particle by a radioactive atom is an electron which is-
\begin{qoptions}
\opt{A}{Emitted from the orbit of the atom.}
\opt{B}{Already present in the nucleus.}
\opt{C}{Produced by the conversion of neutron into proton in the nucleus.}
\opt{D}{Produced by the conversion of proton into neutron in the nucleus.}
\end{qoptions}
\qn{8} The radioactive sources A and B of half lives of 2hr and 4hr respectively, initially contain the same number of radioactive atoms. At the end of 2 hours, their rates of disintegration are in the ratio :-
\begin{qoptions}
\opt{A}{4 : 1}
\opt{B}{2 : 1}
\opt{C}{$\sqrt{2}$ : 1}
\opt{D}{1 : 1}
\end{qoptions}
\qn{9} With the x-axis horizontal and the y-axis vertically upward as positive, the change in the horizontal component of velocity, $\Delta v_x$, and the change in the vertical component of velocity, $\Delta v_y$, of a projectile are related to the time since leaving the barrel, $\Delta t$, as
\begin{qoptions}
\opt{A}{$\Delta v_x = 0;\ \Delta v_y = 0$}
\opt{B}{$\Delta v_x = g\Delta t;\ \Delta v_y = 0$}
\opt{C}{$\Delta v_x = 0;\ \Delta v_y = g\Delta t$}
\opt{D}{$\Delta v_x = 0;\ \Delta v_y = -g\Delta t$}
\end{qoptions}
\qn{10} The net force (in newtons) acting on a body of mass 2kg (at rest) varies with time t (in microseconds) as shown in the figure. AB, BC and CD are straight line segments. Find velocity of the block at 16$\mu$s.\par
\qfig{p2_phy_s1_q10_1}
\begin{qoptions}
\opt{A}{$2.9\times10^{-3}$ m/s}
\opt{B}{$2.9\times10^{3}$ m/s}
\opt{C}{$5\times10^{-3}$ m/s}
\opt{D}{None of these}
\end{qoptions}
\qn{11} A uniform thin rod ABC of mass M and length $\ell$ is placed vertically on a rough horizontal surface. The coefficient of kinetic friction between the rod and the surface is $\mu$. A force F (sufficient for sliding) is applied on the rod at point B at distance $\ell$/3 below centre of the rod as shown in figure. The initial acceleration of point A is\par
\qfig{p2_phy_s1_q11_1}
\begin{qoptions}
\opt{A}{$\mu g-\dfrac{F}{M}$}
\opt{B}{$\dfrac{F}{M}$}
\opt{C}{$4\mu g$}
\opt{D}{$2\mu g-\dfrac{F}{M}$}
\end{qoptions}
\qn{12} A uniform semi circular ring is placed on rough inclined plane as shown. If friction is large for no sliding, then the minimum angle $\theta$ at which toppling occurs is :\par
\qfig{p2_phy_s1_q12_1}
\begin{qoptions}
\opt{A}{$\tan^{-1}2$}
\opt{B}{$\tan^{-1}\dfrac{\pi}{2}$}
\opt{C}{45$^\circ$}
\opt{D}{$\tan^{-1}\dfrac{\pi}{4}$}
\end{qoptions}
\qn{13} As shown in figure, two blocks are connected with a light spring. When spring was at its natural length, velocities are given to them as shown in figure. Choose the wrong alternative.\par
\qfig{p2_phy_s1_q13_1}
\begin{qoptions}
\opt{A}{Velocity of center of mass of the system is 3 m/s (towards right)}
\opt{B}{When spring is maximum compressed velocity of 20 kg block is 3 m/s (towards right)}
\opt{C}{When spring is maximum elongated velocity of 10 kg block is 3 m/s (towards left)}
\opt{D}{At the time of maximum compression both blocks are moving in same direction.}
\end{qoptions}
\qn{14} A square frame of mass m and side $\ell$ is suspended as shown. Find the time period of small oscillations in vertical plane :-\par
\qfig{p2_phy_s1_q14_1}
\begin{qoptions}
\opt{A}{$2\pi\sqrt{\dfrac{\ell}{g}}$}
\opt{B}{$2\pi\sqrt{\dfrac{5\sqrt{2}}{6}\dfrac{\ell}{g}}$}
\opt{C}{$2\pi\sqrt{\dfrac{\ell}{\sqrt{2}g}}$}
\opt{D}{$2\pi\sqrt{\dfrac{5\sqrt{2}}{3}\dfrac{\ell}{g}}$}
\end{qoptions}
\qn{15} Letter F is formed from three identical rods of length $\ell$. Find coordinates of its centre of mass.\par
\qfig{p2_phy_s1_q15_1}
\begin{qoptions}
\opt{A}{$\left(\dfrac{\ell}{5},\dfrac{\ell}{5}\right)$}
\opt{B}{$\left(\dfrac{\ell}{3},\dfrac{\ell}{6}\right)$}
\opt{C}{$\left(\dfrac{\ell}{6},\dfrac{\ell}{3}\right)$}
\opt{D}{$\left(\dfrac{\ell}{2},\dfrac{\ell}{5}\right)$}
\end{qoptions}
\qn{16} Find the work done by friction in slowly moving a block of mass 2kg from A to B on a small hill as shown in figure :-\par
\qfig{p2_phy_s1_q16_1}
\begin{qoptions}
\opt{A}{$-30$ J}
\opt{B}{30 J}
\opt{C}{$-20$ J}
\opt{D}{$-15$ J}
\end{qoptions}
\qn{17} A uniform rod of mass 5m is hinged as shown. Three particles each of mass m are kept as shown. Find acceleration of particle P at t = 0 just after release :-\par
\qfig{p2_phy_s1_q17_1}
\begin{qoptions}
\opt{A}{$\dfrac{18g}{17}$}
\opt{B}{$\dfrac{g}{4}$}
\opt{C}{$\dfrac{g}{16}$}
\opt{D}{$\dfrac{39}{16}g$}
\end{qoptions}
\qn{18} The system is in equilibrium with applied force F as shown. At t = 0 force F is removed. Initial acceleration of 2m block will be :-\par
\qfig{p2_phy_s1_q18_1}
\begin{qoptions}
\opt{A}{$\dfrac{F}{m}$}
\opt{B}{$\dfrac{F}{3m}$}
\opt{C}{$\dfrac{F}{2m}$}
\opt{D}{$\dfrac{2F}{m}$}
\end{qoptions}
\qn{19} A ball collides with a wall and moves parallel to the wall as shown. The impulse acting on the ball during the collision is :-\par
\qfig{p2_phy_s1_q19_1}
\begin{qoptions}
\opt{A}{$mv\left(-\dfrac{4}{5}\hat{i}-\dfrac{1}{10}\hat{j}\right)$}
\opt{B}{$mv\left(-\dfrac{4}{5}\hat{i}-\dfrac{2}{5}\hat{j}\right)$}
\opt{C}{$mv\left(-\dfrac{3}{5}\hat{i}-\dfrac{1}{10}\hat{j}\right)$}
\opt{D}{$mv\left(-\dfrac{3}{5}\hat{i}-\dfrac{2}{5}\hat{j}\right)$}
\end{qoptions}
\qn{20} A sample of H atoms in certain excited state emits 10 different wavelengths in de-excitation to ground state. The maximum wavelength out of these wavelengths is :-
\begin{qoptions}
\opt{A}{$4.05\times10^{-6}$ m}
\opt{B}{$4.05\times10^{-4}$ m}
\opt{C}{$4.5\times10^{-6}$ m}
\opt{D}{$4.5\times10^{-4}$ m}
\end{qoptions}
\sect{SEC-B}
\qn{1} A silver sphere (work function 4.6 eV) is suspended in a vacuum chamber by an insulating thread. Ultraviolet light of wavelength 0.2 $\mu$ m strikes the sphere. The maximum electric potential (in Volts) of the sphere will be (hc = 12400 eV\AA) :-
\qn{2} Rain appears to be falling at 53$^\circ$ with the vertical to a man standing on the ground. When he starts running at 10 km/hr, rain again appears to be falling at 53$^\circ$ with the vertical. Find the speed (in km/hr) of rain wrt ground :-
\qn{3} A projectile is thrown perpendicular to a long inclined plane as shown. The range (in m) of projectile on the incline is :-\par
\qfig{p2_phy_s2_q03_1}
\qn{4} The wavelength of K$_\alpha$ X-ray of an element having atomic number Z = 13 is $\lambda$. The wavelength of K$_\alpha$ X-ray of another element of atomic number Z$'$ is 4$\lambda$. Then Z$'$ is
\qn{5} A sample of radioactive nuclei of a certain element can decay only by $\gamma$-emission and $\beta$-emission. If the half-life for $\gamma$-emission is 24 minutes and that for $\beta$-emission is 36 minutes, the half-life (in minute) for the sample is
\qn{6} A block of metal weighing 2 kg is resting on a frictionless plane. It is struck by a jet releasing water at a rate of 1 kg/s and at a speed of 5 m/s. The initial acceleration (in m/s$^2$) of the block will be (Assume water doesn't rebound) :-\par
\qfig{p2_phy_s2_q06_1}
\qn{7} A ball is thrown with 50 m/s at 30$^\circ$ with the horizontal from a building of height 500 m. The time (in sec) at which velocity vector of the ball becomes perpendicular to the initial velocity is (g = 10 m/s$^2$) :-
\qn{8} A particle starts from rest and performs SHM of amplitude A. Find the ratio of time taken by it from mean to $\dfrac{A\sqrt{3}}{2}$ to time taken by it from $A$ to $\dfrac{A}{2}$ :-
\qn{9} If there will be no relative motion between the two blocks, the friction force (in N) between the blocks will be :-\par
\qfig{p2_phy_s2_q09_1}
\qn{10} A ball is thrown with $40\sqrt{2}$ m/s at 45$^\circ$ with horizontal. Find the radius of curvature (in m) of its path at t = 7 sec :-
\partsec{2}{CHEMISTRY}
\sect{SEC-A}
\qn{1} Which of the following is INCORRECT statement ?
\begin{qoptions}
\opt{A}{In adiabatic expansion of an ideal gas, temperature of system always decrease}
\opt{B}{In isothermal expansion of an ideal gas, temperature of a gas always remains constant}
\opt{C}{In a reversible adiabatic compression of an ideal gas, change in entropy of system is zero}
\opt{D}{In an irreversible adiabatic compression of an ideal gas, change in entropy of system is positive}
\end{qoptions}
\qn{2} Which of the following arrangement does not represent correct increasing order of the property stated against it?
\begin{qoptions}
\opt{A}{Sc Ti V Cr : Highest possible oxidation number}
\opt{B}{Ti V Cr Mn : Melting point}
\opt{C}{Cu Ni Co Fe : Paramagnetic behaviour}
\opt{D}{Mn Fe Co : E$^\circ$ (M$^{2+}$/M) standard reduction potential}
\end{qoptions}
\qn{3} Bone deformities in children is due to deficiency of which vitamins :
\begin{qoptions}
\opt{A}{Ascorbic acid}
\opt{B}{Thiamine}
\opt{C}{Riboflavin}
\opt{D}{Vitamin-D}
\end{qoptions}
\qn{4}\par
\qfig{p2_chem_s1_q04_1}
Which of the following statement is correct.\par
(i) If A \& B are same gases then T$_2$ $>$ T$_1$\par
(ii) If T$_1$ = T$_2$ then A$_{(g)}$ may be SO$_{2(g)}$ \& B$_{(g)}$ may be CH$_{4(g)}$\par
(iii) If B \& C are same gases then T$_3$ $>$ T$_2$\par
(iv) If T$_1$ = T$_2$ = T$_3$ then A$_{(g)}$ may be CH$_4$, B$_{(g)}$ may be SO$_2$ \& C$_{(g)}$ may be SO$_{3(g)}$
\begin{qoptions}
\opt{A}{i, ii \& iii only}
\opt{B}{i, ii, \& iv only}
\opt{C}{All are correct}
\opt{D}{ii \& iii only}
\end{qoptions}
\qn{5} Which of the following is electron precise molecular hydride ?
\begin{qoptions}
\opt{A}{B$_2$H$_6$}
\opt{B}{SiH$_4$}
\opt{C}{NH$_3$}
\opt{D}{BeH$_2$}
\end{qoptions}
\qn{6} In which process hardness of water is removed by the formation of water soluble complex?
\begin{qoptions}
\opt{A}{Calgon's process}
\opt{B}{Zeolite method}
\opt{C}{Clark's method}
\opt{D}{Washing soda method}
\end{qoptions}
\qn{7}\par
\qfig{p2_chem_s1_q07_1}
\begin{qoptions}
\opt{A}{\qfigi{p2_chem_s1_q07_2}}
\opt{B}{\qfigi{p2_chem_s1_q07_3}}
\opt{C}{\qfigi{p2_chem_s1_q07_4}}
\opt{D}{\qfigi{p2_chem_s1_q07_5}}
\end{qoptions}
\qn{8} Identify the monomer in the given polymeric structures :\par
\qfig{p2_chem_s1_q08_1}
\begin{qoptions}
\opt{A}{Propene}
\opt{B}{Styrene}
\opt{C}{Vinylchloride}
\opt{D}{Ethene}
\end{qoptions}
\qn{9} For which of the following reaction entropy, change in the system would be positive (in forward direction)
\begin{qoptions}
\opt{A}{N$_{2(g)}$ + 3H$_{2(g)}$ $\rightarrow$ 2NH$_{3(g)}$}
\opt{B}{3Mg$_{(s)}$ + 2Fe$^{3+}_{(aq)}$ $\rightarrow$ 3Mg$^{2+}_{(aq)}$ + 2Fe$_{(s)}$}
\opt{C}{CH$_{4(g)}$ + 2O$_{2(g)}$ $\rightarrow$ CO$_{2(g)}$ + 2H$_2$O$_{(g)}$}
\opt{D}{C$_{(\text{graphite})}$ $\rightarrow$ C$_{(\text{diamond})}$}
\end{qoptions}
\qn{10} Misch metal consist of :-
\begin{qoptions}
\opt{A}{$\sim$95\% Lanthanoides and $\sim$5\% Fe}
\opt{B}{$\sim$5\% Lanthanoides and 95\% Fe}
\opt{C}{$\sim$95\% Actinides and 5\% Al}
\opt{D}{$\sim$95\% Lanthanoides and $\sim$5\% Al}
\end{qoptions}
\qn{11} Which of the following aldehyde gives positive Iodoform test?
\begin{qoptions}
\opt{A}{\qfigi{p2_chem_s1_q11_1}}
\opt{B}{CH$_3$ $-$ CH = O}
\opt{C}{\qfigi{p2_chem_s1_q11_2}}
\opt{D}{Me$_3$C $-$ CH = O}
\end{qoptions}
\qn{12} Which of the following is non reducing sugar ?
\begin{qoptions}
\opt{A}{Maltose}
\opt{B}{Sucrose}
\opt{C}{Fructose}
\opt{D}{Lactose}
\end{qoptions}
\qn{13} For which of the following process $\Delta$G = 0
\begin{qoptions}
\opt{A}{A$_{(l)}$ $\rightarrow$ A$_{(g)}$ (1 atm, 90$^\circ$C), B.P. of A$_{(l)}$ at 1 atm = 90$^\circ$C}
\opt{B}{H$_2$O$_{(l)}$ $\rightarrow$ H$_2$O$_{(g)}$ (1 atm, 400$^\circ$C)}
\opt{C}{H$_2$O$_{(l)}$ $\rightarrow$ H$_2$O$_{(s)}$ (1.0 atm, $-$10$^\circ$C)}
\opt{D}{A$_{(l)}$ $\rightarrow$ A$_{(g)}$ (2 atm, 90$^\circ$C), B.P. of A$_{(l)}$ at 1 atm = 90$^\circ$C}
\end{qoptions}
\qn{14} Select the CORRECT statement for heavy water.
\begin{qoptions}
\opt{A}{Heavy water has lower boiling point than H$_2$O}
\opt{B}{Heavy water is more effective solvent than ordinary water}
\opt{C}{Heavy water is more associated than ordinary water}
\opt{D}{Ice cube of heavy water floats on H$_2$O(liq.)}
\end{qoptions}
\qn{15} Incorrect statement regarding solution obtained by dissolving alkali metals in liquid ammonia.
\begin{qoptions}
\opt{A}{Paramagnetic nature}
\opt{B}{Show oxidising nature}
\opt{C}{Deep blue due to presence of ammoniated electrons}
\opt{D}{Good conductor}
\end{qoptions}
\qn{16} Which of the following reaction is not correctly matched with their major product?
\begin{qoptions}
\opt{A}{\qfigi{p2_chem_s1_q16_1}}
\opt{B}{\qfigi{p2_chem_s1_q16_2}}
\opt{C}{CH$_3$CH$_2$COCl + (CH$_3$)$_2$Cd $\rightarrow$ CH$_3$COCH$_3$ + CdCl$_2$}
\opt{D}{\qfigi{p2_chem_s1_q16_3}}
\end{qoptions}
\qn{17} Select the correct statement :-\par
(i) Lithium nitrate when heated gives N$_2$ gas\par
(ii) LiF and Li$_2$O are comparatively much less soluble in water than the corresponding compound of alkali metals\par
(iii) LiCl is deliquescent and crystallises as a hydrate, LiCl.2H$_2$O\par
(iv) In KO$_2$, peroxide ion is present which diamagnetic
\begin{qoptions}
\opt{A}{(i), (ii)}
\opt{B}{(ii),(iii)}
\opt{C}{(i),(iv)}
\opt{D}{(iii),(iv)}
\end{qoptions}
\qn{18} Select the correct statement :-
\begin{qoptions}
\opt{A}{MnO$_2$ is used in dry battery cells}
\opt{B}{Lanthanoids react with H$_2$O to formed Ln(OH)$_4$ and H$_2$}
\opt{C}{Lanthanoids react with acids \& do not produce H$_2$}
\opt{D}{In actinoids series common oxidation state is +4}
\end{qoptions}
\qn{19} In which reaction hydrocarbon is not formed as major product :
\begin{qoptions}
\opt{A}{\qfigi{p2_chem_s1_q19_1}}
\opt{B}{\qfigi{p2_chem_s1_q19_2}}
\opt{C}{\qfigi{p2_chem_s1_q19_3}}
\opt{D}{\qfigi{p2_chem_s1_q19_4}}
\end{qoptions}
\qn{20} Correct IUPAC name of given compound is\par
\qfig{p2_chem_s1_q20_1}
\begin{qoptions}
\opt{A}{3-methylcyclohexanecarbaldehyde}
\opt{B}{3-methylcyclohexanal}
\opt{C}{5-Methylcyclohexane carbaldehyde}
\opt{D}{3-Methylcyclohexanol}
\end{qoptions}
\sect{SEC-B}
\qn{1} For the real gases reaction 2A$_{(g)}$ + B$_{(g)}$ $\rightarrow$ 4C$_{(g)}$ ; $\Delta$H = ? In a 10 litre vessel at 500 K the initial pressure is 50 bar. After the complete reaction, volume of gases was found to be 15 litre at a pressure of 80 bar. Find change in enthalpy (in kJ) if change in internal energy is +50 kJ \{1 bar-litre = 100 Joule\}
\qn{2} In acidic medium KMnO$_4$ is decolourised by which of the following specie(s). Give the total sum of the atoms present in them :-\par
(a) H$_2$C$_2$O$_4$ \quad (b) HNO$_3$ \quad (c) Na$_2$S$_2$O$_3$ \quad (d) HNO$_2$
\qn{3} Which of the following is essential amino acid? Give total number of hydrogen atoms present in it.\par
(a) Proline \quad (b) Glutamine \quad (c) Valine \quad (d) Serine
\qn{4} 50 gm H$_2$O(A) is mixed with 70 gm of H$_2$O(B) in a rigid adiabatic vessel. The initial temperature of A \& B was 46$^\circ$ \& 70$^\circ$C. Find $\Delta$H$_A$ [change in enthalpy of water H$_2$O(A)] (in calories) specific heat of water = 1 cal/gm/$^\circ$c
\qn{5}\par\noindent{\footnotesize\begin{tabular}{@{}l@{\ }l@{}}
HCl + BOH $\rightarrow$ BCl + H$_2$O & $\Delta$H = $-$57 kJ\\
CH$_3$COOH + BOH $\rightarrow$ CH$_3$COOB + H$_2$O & $\Delta$H = $-$50 kJ\\
HCl + DOH $\rightarrow$ DCl + H$_2$O & $\Delta$H = $-$45 kJ
\end{tabular}}\par
then $\Delta$H (in kJ) for CH$_3$COOH + DOH $\rightarrow$ CH$_3$COOD + H$_2$O, is equal to :-
\qn{6} In one of the following reaction H$_2$O$_2$ shows reducing nature. Find the n-factor of the species getting reduced ?\par
(1) CrO$_4^{2-}$ + 2H$_2$O$_2$ $\xrightarrow{\text{H}^+}$ CrO$_5$ + 3H$_2$O\par
(2) MnO$_2$ + H$_2$O$_2$ + H$_2$SO$_4$ $\longrightarrow$ MnSO$_4$ + 2H$_2$O + O$_2$\par
(3) H$_2$S + H$_2$O$_2$ $\longrightarrow$ S$\downarrow$ + 2H$_2$O\par
(4) Mn$^{+2}$ + H$_2$O$_2$ $\xrightarrow{\text{OH}^-}$ MnO$_2$ + 2H$_2$O
\qn{7} The molecular speed in a sample of 100 molecules are distributed as\par
\noindent\begin{tabular}{@{}l|cccccc@{}}
Number of molecules & 10 & 20 & 40 & 15 & 10 & 5\\
Speed m/sec. & 60 & 80 & 100 & 120 & 140 & 160
\end{tabular}\par
Give the sum of V$_{\text{mp}}$ \& V$_{\text{avg.}}$ in m/sec
\qn{8}\par
\qfig{p2_chem_s2_q08_1}
Give the total number of atoms present in final product `B'.
\qn{9} The critical point of CO$_2$ is 304 K and 73 atm. Under which condition is carbon dioxide liquid.\par
I) 305 K, 70 atm \quad II) 303 K, 75 atm\par
III) 303 K, 73 atm \quad IV) 305 K, 80 atm\par
Give the ratio $\dfrac{\text{sum of temperature (in K) of all}}{\text{sum of pressure (in atm) of all}}$ these conditions
\qn{10} At what temperature (in K) is the following reaction at equilibrium, when all substances are at standard pressure? Assuming that entropies and enthalpies of reaction do not vary with temperature.\par
PCl$_3$ + Cl$_{2(g)}$ $\rightleftharpoons$ PCl$_{5(g)}$\par
\noindent\begin{tabular}{@{}l@{\quad}c@{\quad}c@{}}
 & $\Delta$H$_f^0$ kJ/mole & S$^\circ$J/mole/kelwing\\
PCl$_{3(g)}$ & $-$290 & 310\\
Cl$_{2(g)}$ & 0 & 220\\
PCl$_{5(g)}$ & $-$375 & 360
\end{tabular}
\partsec{3}{MATHEMATICS}
\sect{SEC-A}
\qn{1} The value of $\displaystyle\int_0^{10\pi}[\sin x+\cos x]\,dx$ is (where [.] denotes greatest integer function)
\begin{qoptions}
\opt{A}{$-5\pi$}
\opt{B}{$-10\pi$}
\opt{C}{$-\dfrac{5\pi}{2}$}
\opt{D}{$-\dfrac{15\pi}{2}$}
\end{qoptions}
\qn{2} Let $f:\mathbb{R}\to\mathbb{R}$ be a continuous function. If $f(x)+f(x+1)+f(x+2)+f(x+3)=2021,\ \forall\ x\in\mathbb{R}$ and $\displaystyle\int_0^4 f(x)\,dx=A$, then the value of $\displaystyle\int_0^{100}f(x)\,dx$ is equal to -
\begin{qoptions}
\opt{A}{A}
\opt{B}{20A}
\opt{C}{25A}
\opt{D}{50A}
\end{qoptions}
\qn{3} If $f(x)=2(g(x))^3-3(g(x))^2+3g(x)-2x+\cos x$, where $g(x)$ is a differentiable function then which of the following is always true-
\begin{qoptions}
\opt{A}{$f(x)$ is increasing if $g(x)$ is decreasing}
\opt{B}{$f(x)$ is decreasing if $g(x)$ is increasing}
\opt{C}{$f(x)$ is increasing if $g(x)$ is increasing}
\opt{D}{$f(x)$ is decreasing if $g(x)$ is decreasing}
\end{qoptions}
\qn{4} Angle between the tangents drawn at origin to the two curves $x^5+4x^3y^2+3x+4y=0$ and $x^3-27y^3+5(4x-3y)^2=0$ is
\begin{qoptions}
\opt{A}{$\tan^{-1}1+\tan^{-1}2+\tan^{-1}3$}
\opt{B}{$\cot^{-1}1+\cot^{-1}2+\cot^{-1}3$}
\opt{C}{is the angle between the lines $11y=2x$ and $x+2y-0=0$}
\opt{D}{$2\cot^{-1}1+\cot^{-1}2$}
\end{qoptions}
\qn{5} The point on the curve $y^3=4x^2+12y$ at which the tangent is horizontal is
\begin{qoptions}
\opt{A}{(2,2)}
\opt{B}{($-$2,$-$2)}
\opt{C}{(0,0)}
\opt{D}{(1,1)}
\end{qoptions}
\qn{6} Equation of a straight line with gradient 1 which touches curve $y=\dfrac{x}{1+|x|}$ is
\begin{qoptions}
\opt{A}{$y=x$}
\opt{B}{$y=x+1$}
\opt{C}{$y=x-1$}
\opt{D}{$y=x+2$}
\end{qoptions}
\qn{7} If A is a symmetric matrix ($|A|=-1$) and $AA^T=I$, then identify the incorrect option-
\begin{qoptions}
\opt{A}{$A^T=A^{-1}$}
\opt{B}{$A^T+\mathrm{adj}A=0$}
\opt{C}{$A^2=A^4$}
\opt{D}{$A^2=A^3$}
\end{qoptions}
\qn{8} The function $f(x)=\displaystyle\int_{-2}^{x}(e^{-t}-1)(t^2-1)^5(t+2)^9(t-3)^7\,dt$ has a local minimum at `x' equal to
\begin{qoptions}
\opt{A}{$-2$}
\opt{B}{0}
\opt{C}{3}
\opt{D}{$-1$}
\end{qoptions}
\qn{9} If $f(x)$ is an even function such that $f(10-x)=f(x)$ and $\displaystyle\int_0^5 f(x)\,dx=7$. Then the value of $\displaystyle\int_0^{10}f(x)\,dx$ is
\begin{qoptions}
\opt{A}{10}
\opt{B}{12}
\opt{C}{14}
\opt{D}{21}
\end{qoptions}
\qn{10} If $f(x)=\displaystyle\int\left(2\sin(2\ln x)-\cos(2\ln x)\right)dx$ and $f\!\left(e^{\frac{\pi}{2}}\right)=0$. Then the value of $f(e^{\pi/4})$ is
\begin{qoptions}
\opt{A}{$-e^{\pi/2}$}
\opt{B}{$e^{\pi/4}\left(1-2e^{\pi/4}\right)$}
\opt{C}{$e^{\pi/4}\left(2e^{\pi/4}-1\right)$}
\opt{D}{$e^{\pi/4}\left(2e^{\pi/2}-1\right)$}
\end{qoptions}
\qn{11} If length of subnormal to the curve $x\left(\dfrac{y}{k}\right)^n=k$ is always constant then `n' is equal to
\begin{qoptions}
\opt{A}{2}
\opt{B}{1}
\opt{C}{$-1$}
\opt{D}{$-2$}
\end{qoptions}
\qn{12} Let A and B are orthogonal matrices of order 3 $\times$ 3 such that $X=AB^T$ and $Y=BA^T$, then the value of tr(XY) is equal to -
\begin{qoptions}
\opt{A}{1}
\opt{B}{2}
\opt{C}{3}
\opt{D}{4}
\end{qoptions}
\qn{13} Number of point(s) of inflection of $f(x)=x^4+6x^3-35x^2+6x-3$ is -
\begin{qoptions}
\opt{A}{0}
\opt{B}{1}
\opt{C}{2}
\opt{D}{4}
\end{qoptions}
\qn{14} Solution of differential equation $\left(\sqrt{1-x^2y^2}-y\right)dx=x\,dy$ cannot be
\begin{qoptions}
\opt{A}{$xy+\cos x=0$}
\opt{B}{$xy=\sin x$}
\opt{C}{$xy+\sin x=0$}
\opt{D}{$xy=\tan x$}
\end{qoptions}
\qn{15} The curve which satisfy differential equation $(x^2y-y)dx+(xy^2-x)dy=0$ is-
\begin{qoptions}
\opt{A}{$\pi x^2y^2=e^{x^2+y^2}$}
\opt{B}{$\pi x^2y^2=e^{x+y}$}
\opt{C}{$exy=e^{x+y}$}
\opt{D}{$exy=e^{x^2+y^2}$}
\end{qoptions}
\qn{16} A function $y=f(x)$ satisfies the differential equation $\dfrac{dy}{dx}+x^2y+2x=0,\ f(1)=1$, then the value of $f'(1)$ is-
\begin{qoptions}
\opt{A}{1}
\opt{B}{$-1$}
\opt{C}{3}
\opt{D}{$-3$}
\end{qoptions}
\qn{17} Equation of tangent drawn to the curve $f(x)=x^3\dfrac{(2^x-1)}{(2^x+1)}+\sin^2x\cos^6x-1$ at its point of intersection with y-axis is-
\begin{qoptions}
\opt{A}{$y+1=0$}
\opt{B}{$x+3y=10$}
\opt{C}{$x+y=0$}
\opt{D}{$x=-1$}
\end{qoptions}
\qn{18} Consider a system of equations $x-y+z=-1$, $\lambda x+3y+2z=3$ and $3x+\lambda y+z=2$, if system have a unique solution then $\lambda$ \textbf{cannot} be-
\begin{qoptions}
\opt{A}{3}
\opt{B}{$-3$}
\opt{C}{$-4$}
\opt{D}{2}
\end{qoptions}
\qn{19} If $A=\begin{bmatrix}3&2&1\\1&2&1\\0&1&2\end{bmatrix}$ also, $|\mathrm{adj}\,\lambda A|=\lambda^{a}.a^{\lambda}$ (where $a,\lambda\in N$). Then the value of `$\alpha+\lambda$', is-
\begin{qoptions}
\opt{A}{6}
\opt{B}{8}
\opt{C}{10}
\opt{D}{12}
\end{qoptions}
\qn{20} If $A=\begin{pmatrix}2&7\\1&3\end{pmatrix}$, then the value of $\det(A^{2019}-3A^{2018})$ is equal to-
\begin{qoptions}
\opt{A}{$-7$}
\opt{B}{7}
\opt{C}{5}
\opt{D}{$-5$}
\end{qoptions}
\sect{SEC-B}
\qn{1} If $\displaystyle\int e^x\left(\dfrac{\tan x}{1+\tan x}+\dfrac{1}{1+\sin 2x}-\dfrac{1}{2}\right)dx=e^xf(x)+c$, where c is integration constant. Then $\left[5f\!\left(\dfrac{5\pi}{8}\right)\right]$ is (where [.] denotes greatest integer function)
\qn{2} If $f(x)=$ {\scriptsize$\begin{vmatrix}3x^2-6x+5&6x^2+6x-3&3x^2+3x-2\\2x-2&4x+2&2x+1\\1&2&1\end{vmatrix}$}, then the value of $\dfrac{f(1)+f(3)}{f(2)}$
\qn{3} Let $f(x)$ be a thrice differentiable function such that it intersects the curve $g(x)=2x^3+7x^2-5x+1$ at exactly 7 distinct points, then minimum number of roots of $f'''(x)=12$ is
\qn{4} The area bounded by $y^2=4x$ and $x^2=4y$ is
\qn{5} The area between the curve $y=2x^4-x^2$, the x-axis and the ordinates of two minima of the curve is of the form $\dfrac{a}{b}$ (where a and b are co-prime). Then the value of $(b-10a)$ is-
\qn{6} If $f(x)=2000x^3-3000x^2-36000x+2021$ has local maximum at $x=p$ and local minimum at $x=q$, then the value of $\left|\dfrac{q}{p}\right|$ is
\qn{7} Area (in square units) bounded by the curves $f(x)=(\cos x+\sec x)^2$ and $g(x)=(\sin x+\mathrm{cosec}\,x)^2$ between $x=\dfrac{\pi}{4}$ and $x=\dfrac{\pi}{3}$ is of the form $\dfrac{19\sqrt{b}-c}{d}$, s.t. H.C.F of b,c and d is 3, then the value of $\dfrac{b+c}{d}$ is-
\qn{8} The area bounded by pair of straight lines $x^2+y^2-2xy-5x+5y+6=0$ and the co-ordinate axes is-
\qn{9} Minimum distance between the curves $C_1: x^2+y^2=8$ and $C_2: xy=9$ is-
\qn{10} If $f(x)=x^3+3ax^2+(11a-6)x-2$ is increasing function for all $x\in\mathbb{R}$, then number of possible integral value(s) of `a' is/are-
\end{multicols}
\end{document}
